% ============================================================
%  Dublin — reusable components
%  Loaded after dublin-theme.tex.
% ============================================================

\tcbset{
  dgbase/.style={
    enhanced, sharp corners, boxrule=0pt, frame hidden,
    left=1.1em, right=1.1em, top=0.75em, bottom=0.75em,
    fontupper=\footnotesize, halign=flush left
  }
}

% ---------- callout boxes ----------
% Five boxes, one shape: a flat wash panel with a coloured bar down the left.
% Every one takes an optional title, so there is a single thing to remember.
% The bar colour is the only difference between them.
%
%   \begin{takeaway}[What this changes] ... \end{takeaway}
%   \begin{result} ... \end{result}
%
% From markdown the Lua filter turns
%   ::: {.takeaway title="What this changes"}
% into the same thing, and drops the title if there is none.

% One header format for every box: the title on its own line, bold, in the
% box's own colour, set as typed. It used to be \MakeUppercase and a size
% down, which made a takeaway look like a different component from a defn.
\newcommand{\dgboxtitle}[3]{%
  % #1 title colour, #2 title text, #3 sets whether a title was given
  \ifx\relax#2\relax\else
    {\color{#1}\bfseries #2\par\vspace{0.15em}}%
  \fi
  \color{dgInk}\mdseries#3%
}

\newtcolorbox{dgcallout}[2][dgGreen]{%
  dgbase, colback=dgWash, borderline west={3pt}{0pt}{#1}, #2
}

% takeaway -- the one thing they should leave with
\newenvironment{takeaway}[1][]{%
  \begin{dgcallout}[dgGreen]{}\dgboxtitle{dgGreen}{#1}{}%
}{\end{dgcallout}}

% result -- a finding, stated flatly
\newenvironment{result}[1][]{%
  \begin{dgcallout}[dgBlue]{}\dgboxtitle{dgBlue}{#1}{}%
}{\end{dgcallout}}

% caveat -- a limitation or an assumption
\newenvironment{caveat}[1][]{%
  \begin{dgcallout}[dgMuted]{}\dgboxtitle{dgMuted}{#1}{}%
}{\end{dgcallout}}

% defn -- a definition, for teaching. Title runs inline with the body.
\newenvironment{defn}[1][]{%
  \begin{dgcallout}[dgNavy]{}\dgboxtitle{dgNavy}{#1}{}%
}{\end{dgcallout}}

% worked -- a worked example, for teaching
\newenvironment{worked}[1][]{%
  \begin{dgcallout}[dgBlue]{}\dgboxtitle{dgBlue}{#1}{\footnotesize}%
}{\end{dgcallout}}

% tablenotes -- how to read the exhibit above. Same panel and bar as the
% rest, one size smaller, with the label running inline as in defn. Set
% from markdown as ::: {.tablenotes}, via \dgnotes in the Lua filter.
%
% \long, because the filter passes the note through as pandoc blocks
% rather than a flattened string, so the argument can hold a paragraph
% break.
\newtcolorbox{dgnotesbox}{%
  dgbase, colback=dgWash,
  borderline west={3pt}{0pt}{dgMuted}, top=0.55em, bottom=0.55em
}
\makeatletter
\newcommand{\dgnoteslabel}{Notes}
\long\def\dgnotes{\@ifnextchar[\dgnotes@opt\dgnotes@plain}
\long\def\dgnotes@plain#1{\dgnotes@do{\dgnoteslabel}{#1}}
\def\dgnotes@opt[#1]{\dgnotes@rest{#1}}
\long\def\dgnotes@rest#1#2{\dgnotes@do{#1}{#2}}
\long\def\dgnotes@do#1#2{%
  \par\vspace{0.35em}%
  \begin{dgnotesbox}%
    {\color{dgMuted}\bfseries #1\par\vspace{0.15em}}%
    \color{dgInk}\mdseries#2%
  \end{dgnotesbox}%
}
\makeatother

% Kept so older decks still build; \begin{takeaway}[Title] is the way now.
\newenvironment{takeawaytitled}[1]{\begin{takeaway}[#1]}{\end{takeaway}}

% ---------- numbered contributions ----------
% \contribution{1}{Heading}{Body text.}
%
% \long, because the Lua filter now passes the heading and the body through
% as pandoc blocks rather than as flattened LaTeX strings — which is what
% lets a citation inside a .contributions item reach citeproc. A pandoc
% block carries a paragraph break, so a non-\long macro would fail on it.
\long\def\contribution#1#2#3{%
  \par\vspace{0.5em}
  \noindent
  \begin{minipage}[t]{1.6em}\raggedright
    {\color{dgGreen}\bfseries\large #1}
  \end{minipage}%
  \begin{minipage}[t]{\dimexpr\linewidth-1.9em\relax}\raggedright
    % \par rather than \\ : the heading now arrives as a pandoc block and
    % carries its own paragraph break, and \\ after a \par is "There's no
    % line here to end". \par-then-\vspace is correct whether or not the
    % argument ends a paragraph itself.
    {\color{dgNavy}\bfseries #2\par}
    \vspace{0.15em}
    {\small #3}
  \end{minipage}
  \par\vspace{0.35em}
}

% ---------- headline number ----------
% \statline{0.22pp}{lower annual income growth at median exposure}
\newcommand{\statline}[2]{%
  \begin{minipage}[t]{\linewidth}\raggedright
    {\color{dgNavy}\bfseries\fontsize{30}{32}\selectfont #1}\\[0.25em]
    {\color{dgMuted}\small #2}
  \end{minipage}%
}
% Two across: \statrowii{a}{la}{b}{lb}
\newcommand{\statrowii}[4]{%
  \noindent
  \begin{minipage}[t]{0.48\linewidth}\statline{#1}{#2}\end{minipage}\hfill
  \begin{minipage}[t]{0.48\linewidth}\statline{#3}{#4}\end{minipage}%
}
% Three across: \statrow{a}{la}{b}{lb}{c}{lc}
\newcommand{\statrow}[6]{%
  \noindent
  \begin{minipage}[t]{0.32\linewidth}\statline{#1}{#2}\end{minipage}\hfill
  \begin{minipage}[t]{0.32\linewidth}\statline{#3}{#4}\end{minipage}\hfill
  \begin{minipage}[t]{0.32\linewidth}\statline{#5}{#6}\end{minipage}%
}

% ---------- appendix navigation ----------
% Three levels, so a deck can signal how far off-path a jump goes.
%   \jumpto{label}{Text}     level 1, green          — the one you expect
%                                                       to be asked
%   \jumpto[2]{label}{Text}  level 2, blue on white  — the supporting slide
%   \jumpto[3]{label}{Text}  level 3, navy           — background detail
%   \backto{label}{Back}     level 3                 — the return
%
% Three levels, and the return is simply the third of them: a way back is the
% quietest thing in the row. Level 1 is the loudest because the row exists to
% mark the key link, not to list everything behind the talk. Green goes on
% the slide you wrote the appendix for; the quieter levels carry the rest.
% Mark the destination with \hypertarget{label}{} on that frame.
\newcommand{\dgpill}[4]{%
  \hyperlink{#1}{%
    \tikz[baseline=(b.base)]{
      \node[fill=#3, text=#4, font=\bfseries\scriptsize,
            inner xsep=0.6em, inner ysep=0.4em] (b) {\MakeUppercase{#2}};}}%
}
% One level scale, shared by jumps and returns, so a back button is the
% same object in a different colour rather than a different object.
%   White type throughout. 0 and 3 navy, 1 green, 2 blue.
\newcommand{\dglevel}[3]{%
  \ifcase#1\relax
        \dgpill{#2}{#3}{dgNavy}{white}%   % 0, kept as an alias of level 3
    \or \dgpill{#2}{#3}{dgGreen}{white}%
    \or \dgpill{#2}{#3}{dgBlue}{white}%
    \or \dgpill{#2}{#3}{dgNavy}{white}%
  \fi
}
\newcommand{\jumpto}[3][1]{\dglevel{#1}{#2}{#3}}
% The return is level 3: a way back is the quietest thing in the row.
\newcommand{\backto}[3][3]{\dglevel{#1}{#2}{#3}}

% A row of jumps, pushed to the foot of the frame:
%     \navrow{\jumpto{rob}{Robustness} \jumpto[3]{data}{Data}}
\newcommand{\navrow}[1]{\vfill\noindent #1\vspace{0.6em}}

% ---------- contents ----------
% A navigable contents slide, for a lecture where students need to see the
% shape of the hour. Entries are live links: clicking one jumps to that
% section. Content-only, so it goes inside an ordinary "##" frame.
% Entries are live links -- clicking one jumps to that section -- and they
% inherit the same template the section dividers use, so the contents slide
% and the dividers read as one list rather than two.
% Keep this in the ordinary vertical list. A minipage here is as tall as the
% text area allows and pushes the footer off the slide, and so does changing
% \rightskip inside the group -- both were tried.
% The first \tableofcontents in a document leaves a stray skip behind when it
% reads the .toc file, which is enough to push the footer off the slide. A
% box of fixed height bounds it, so the contents slide sits the same whether
% it is the first one in the deck or the fifth.
\newcommand{\toccontent}{%
  \begin{minipage}[t][0.70\textheight][t]{\linewidth}
    \raggedright
    \tableofcontents[hideallsubsections]
  \end{minipage}%
}

% ---------- handouts ----------
% For teaching: print several slides to a page for students. Put one of these
% in header-includes, and add `handout` to classoption so overlays collapse.
%     classoption: [t, handout]
\newcommand{\handouttwoup}{%
  \pgfpagesuselayout{2 on 1}[a4paper, border shrink=5mm]%
}
\newcommand{\handoutfourup}{%
  \pgfpagesuselayout{4 on 1}[a4paper, landscape, border shrink=5mm]%
}

% ---------- references ----------
% Quarto renders citations with citeproc into a CSLReferences environment.
% For APA 7, set in the front matter:
%     bibliography: refs.bib
%     csl: https://www.zotero.org/styles/apa
% References run long, so give the frame [allowframebreaks].
% The list starts hard against the rule under the frame title otherwise.
\AtBeginEnvironment{CSLReferences}{%
  \vspace{0.5em}\footnotesize\setlength{\parskip}{0.45em}%
}

% ---------- quiet source line under an exhibit ----------
\newcommand{\source}[1]{%
  \par\vspace{0.35em}{\color{dgMuted}\scriptsize #1\par}%
}

% ---------- title-slide badge ----------
% \badge{Job Market Paper} — place inside \title{} or on the title slide
\newcommand{\badge}[1]{%
  \tikz[baseline=(b.base)]{
    \node[fill=dgGreen, text=dgNavy, font=\bfseries\scriptsize,
          inner xsep=0.55em, inner ysep=0.35em] (b) {\MakeUppercase{#1}};
  }%
}

% ---------- braced terms in equations ----------
% \dgbrace{i_t - \pi_t - r^*}{real interest rate, relative to the natural rate}
%
% A plain \underbrace sets the label at full width, so a label longer than
% the term it sits under pads the whole equation out and the spacing goes
% wrong. This sets the label in a zero-width box, centred: the equation
% spaces as if the label were not there, and the label is free to overhang.
%
% One brace per term, and never nested — a brace inside a brace is harder to
% read than the algebra it explains. Use it at a composite term's FIRST
% appearance, to say what the combination means.
\newcommand{\dgbrace}[2]{%
  \underbrace{#1}_{\makebox[0pt]{\normalfont\footnotesize #2}}%
}

% \dgbraceover{-\alpha(r_t - r^*) + \varepsilon^y_t}{the short-run deviation}
%
% The same thing above the term instead of below it, for when a WHOLE side of
% an equation is being named — typically to say that a block just written is
% the same object as something defined earlier. Use it sparingly: one per
% equation, and never in the same equation as a \dgbrace on the same term.
\newcommand{\dgbraceover}[2]{%
  \overbrace{#1}^{\makebox[0pt]{\normalfont\footnotesize #2}}%
}

% ---------- struck-out term in an equation ----------
% \dgstrike{(X - M)}   or   \dgstrike[dgBlue]{...}
%
% Rules a coloured line through a term, for showing an assumption being
% imposed — "we close the economy, so net exports go". Works in text and in
% maths, and keeps the term's natural width so the equation does not reflow.
%
% Deliberately NOT \cancel: that needs the cancel package, which is not in
% this theme and cannot be assumed present on a MiKTeX install with
% automatic package installation switched off. tikz is already loaded.
\newcommand{\dgstrike}[2][red]{%
  \tikz[baseline=(dgst.base)]{%
    \node[inner sep=0pt, outer sep=0pt] (dgst) {$#2$};
    \draw[#1, line width=0.9pt] (dgst.south west) -- (dgst.north east);%
  }%
}

% ---------- substitution: what is being replaced, and where it lands ----------
% \dgsub{y_t - y_t^{*}}        the term, in the alert colour
% \dgsubbg{-\alpha(\beta_\pi - 1)(\pi_t - \pi^*) + \varepsilon^y_t}
%                              the block that has just landed, tinted
%
% A substitution is the step students lose. Two equations are on the slide,
% one of them is about to go inside the other, and nothing on the page says
% WHICH PART goes WHERE. The rule is:
%
%   1. On the frame that sets the two equations up, mark the SAME object in
%      both with \dgsub -- the output gap on the left of the IS-MP curve and
%      the output gap on the right of the Phillips curve are one object, so
%      they are one colour.
%   2. On the frame that does the substitution, mark the block that has just
%      arrived with \dgsubbg, in its new surroundings. That is the "where it
%      went", and it is the half students cannot reconstruct on their own.
%   3. Drop the colour the moment the algebra moves on. It marks one step,
%      not a variable: a term that stays red for three slides has stopped
%      meaning anything.
%
% One substitution per pair of frames. Two colours on one slide is two
% things to track, which is the problem this solves, not the solution.
\newcommand{\dgsub}[1]{{\color{dgAlert}#1}}

% The landed block. tikz rather than \colorbox because \colorbox does not
% behave in display maths, and tikz is already loaded for \dgstrike.
\newcommand{\dgsubbg}[1]{%
  \tikz[baseline=(dgsb.base)]{%
    \node[fill=dgAlertBg, rounded corners=1.5pt, inner xsep=2.5pt,
          inner ysep=1.5pt, outer sep=0pt] (dgsb)
      {$\color{dgAlert}#1$};%
  }%
}

% ---------- figures ----------
% Content-only helpers, so they work inside a Quarto "##" frame.
%   \fig{path}{caption}                 one exhibit, sized to the frame
%   \figpair{pathA}{capA}{pathB}{capB}  two side by side
%
% The caption names the exhibit and autonumbers. Source and detail go
% underneath in \dgnotes{} — do not fold them into the caption.
% \tabcaption{...} is the table equivalent; put it above the table, as in
% a paper, with \dgnotes{} below.
\newcommand{\figcaption}[1]{%
  \refstepcounter{figure}%
  {\raggedright\footnotesize\color{dgInk}%
   \textbf{Figure \thefigure:}\ #1\par}%
}
\newcommand{\tabcaption}[1]{%
  \refstepcounter{table}%
  {\raggedright\footnotesize\color{dgInk}%
   \textbf{Table \thetable:}\ #1\par\vspace{0.15em}}%
}
% Height is 0.42\textheight, not 0.56. At 0.56 the exhibit plus a one-line
% caption reaches y = 197 on a 16:9 frame at 11pt, which leaves 47pt before
% the footer -- less than the 53pt a two-line callout needs, and nothing at
% all if the caption runs to two lines. Every figure frame in the decks
% overran for that reason. 0.48 buys about 12pt, so the frame carries the
% exhibit, its caption AND a two-line reading, with the fuller version on the
% next frame or in ::: {.notes}.
\newcommand{\fig}[2]{%
  {\centering
   \includegraphics[width=\linewidth, height=0.42\textheight, keepaspectratio]{#1}\par}%
  \ifx\relax#2\relax\else\vspace{0.45em}\figcaption{#2}\fi
}
% Two figures side by side, each numbered in its own right: a pair is two
% exhibits, not one exhibit in two panels.
\newcommand{\figpair}[4]{%
  \noindent
  \begin{minipage}[t]{0.48\linewidth}
    {\centering\includegraphics[width=\linewidth, height=0.44\textheight,
       keepaspectratio]{#1}\par}%
    \vspace{0.35em}\figcaption{#2}%
  \end{minipage}\hfill
  \begin{minipage}[t]{0.48\linewidth}
    {\centering\includegraphics[width=\linewidth, height=0.44\textheight,
       keepaspectratio]{#3}\par}%
    \vspace{0.35em}\figcaption{#4}%
  \end{minipage}%
}

% ---------- full-bleed figure frame ----------
% \figureframe{Frame title}{path/to/fig.pdf}{Source note.}
% \figtriple{a.png}{caption A}{b.png}{caption B}{c.png}{caption C}
%
% Three panels across, for a comparison that needs a third term - weak,
% medium, strong; low, normal, high. Two panels are \figpair; three are this.
% At 0.32\linewidth each panel is about 45mm on a 16:9 slide, which is small
% but readable for a scatter or a density. Anything that needs to be read off
% an axis precisely wants \figpair and a dropped middle term instead.
\newcommand{\figtriple}[6]{%
  \noindent
  \begin{minipage}[t]{0.32\linewidth}
    {\centering\includegraphics[width=\linewidth, height=0.44\textheight,
       keepaspectratio]{#1}\par}%
    \vspace{0.35em}\figcaption{#2}%
  \end{minipage}\hfill
  \begin{minipage}[t]{0.32\linewidth}
    {\centering\includegraphics[width=\linewidth, height=0.44\textheight,
       keepaspectratio]{#3}\par}%
    \vspace{0.35em}\figcaption{#4}%
  \end{minipage}\hfill
  \begin{minipage}[t]{0.32\linewidth}
    {\centering\includegraphics[width=\linewidth, height=0.44\textheight,
       keepaspectratio]{#5}\par}%
    \vspace{0.35em}\figcaption{#6}%
  \end{minipage}%
}

\newcommand{\figureframe}[3]{%
  \begin{frame}{#1}
    \centering
    \includegraphics[width=\linewidth, height=0.72\textheight, keepaspectratio]{#2}
    \ifx\relax#3\relax\else\source{#3}\fi
  \end{frame}%
}

% ---------- closing frame ----------
% \thanksframe{sam.deegan@ucdconnect.ie}{sam-deegan.com}
\newcommand{\thanksframe}[2]{%
  \begingroup
  \begin{frame}[plain,noframenumbering]
    \thankscontent{#1}{#2}
  \end{frame}
  \endgroup
}
% ---------- appendix divider ----------
\newcommand{\appendixframe}{%
  \appendix
  \dgappendixstart
  \begingroup
  \setbeamercolor{background canvas}{bg=dgWash}
  \begin{frame}[plain,noframenumbering]
    \vfill
    \begin{minipage}{\dimexpr\paperwidth-2.6cm\relax}
      \raggedright
      {\color{dgBlue}\bfseries\small\MakeUppercase{Backup}\par}
      \vspace{0.35em}
      {\color{dgNavy}\bfseries\LARGE Appendix\par}
      \vspace{0.7em}{\color{dgGreen}\rule{2.2cm}{2pt}}
    \end{minipage}
    \vfill
  \end{frame}
  \endgroup
}

% ============================================================
%  Frame-safe variants, for Quarto
%
%  \thanksframe and \appendixframe above create their own frame, which
%  suits plain LaTeX. Under Quarto a "##" heading has already opened the
%  frame, so use these inside a `## {.plain}` slide instead.
% ============================================================

% Paint the whole page a colour from inside an open frame.
\newcommand{\fullbleed}[1]{%
  \begin{tikzpicture}[remember picture, overlay]
    \fill[#1] (current page.south west) rectangle (current page.north east);
  \end{tikzpicture}%
}

\makeatletter
\newcommand{\thankscontent}[2]{%
  \begin{tikzpicture}[remember picture, overlay]\dg@coverbg{0.12}\end{tikzpicture}%
  \vfill
  \begin{minipage}{\dimexpr\paperwidth-2.6cm\relax}
    \raggedright
    {\color{white}\bfseries\LARGE Thank you\par}
    \vspace{0.8em}{\color{dgGreen}\rule{2.6cm}{2pt}\par}
    \vspace{0.8em}
    {\color{white}\normalsize #1\par}
    {\color{dgRule}\normalsize #2\par}
  \end{minipage}
  \vfill
  \ifx\dg@thanksqr\@empty\else
    \begin{tikzpicture}[remember picture, overlay]
      \node[anchor=south east, fill=white, inner sep=5pt, rounded corners=2pt]
        at ($(current page.south east)+(-1.4cm,1.4cm)$)
        {\includegraphics[width=2.4cm]{\dg@thanksqr}};
    \end{tikzpicture}%
  \fi
}
\makeatother

\newcommand{\appendixcontent}{%
  \dgappendixstart
  \fullbleed{dgWash}%
  \vfill
  \begin{minipage}{\dimexpr\paperwidth-2.6cm\relax}
    \raggedright
    {\color{dgBlue}\bfseries\small\MakeUppercase{Backup}\par}
    \vspace{0.35em}{\color{dgNavy}\bfseries\LARGE Appendix\par}
    \vspace{0.7em}{\color{dgGreen}\rule{2.2cm}{2pt}}
  \end{minipage}
  \vfill
}
